% !TeX root = ../../Book5.tex
% 纲要标题：### 13.4 经典例子
% 纲要位置：工作记录/计划纲要.md，第 4505 行。

\section{经典例子}
\label{b5:sec:1304}

% 纲要范围：
%
% * sl_n、so_n 与 sp_2n
% * 维数、中心与单性的范围
% * 小秩同构及例外情形
% * 用括号传播与特征标投影解释单性证明，并以外幂配对展开四维低秩同构
%
% ---
%

经典矩阵 Lie 代数把上述判据变成可以直接计算的例子。维数、中心和单性可以由矩阵单位、型及伴随作用直接证明。后面的根系将解释这些计算为什么分成四个系列。

\subsection{\texorpdfstring{\(\mathfrak{sl}_n\)、\(\mathfrak{so}_n\) 与 \(\mathfrak{sp}_{2n}\)}{sln、son 与 sp2n}}
\label{b5:subsec:1304:01}

对复向量空间，非退化对称型可选基写成单位矩阵，非退化交替型可写成标准辛矩阵。因此采用下列矩阵实现：
\[
\begin{aligned}
\mathfrak{sl}_n(\mathbb C)&=\bigl\{\,X\in M_n(\mathbb C)\mid\operatorname{tr}X=0\,\bigr\},\\
\mathfrak{so}_n(\mathbb C)&=\bigl\{\,X\in M_n(\mathbb C)\mid X^{\mathsf T}+X=0\,\bigr\},\\
\mathfrak{sp}_{2n}(\mathbb C)&=\bigl\{\,X\in M_{2n}(\mathbb C)\mid
X^{\mathsf T}J+JX=0\,\bigr\},\qquad
J=\begin{pmatrix}0&I_n\\-I_n&0\end{pmatrix}.
\end{aligned}
\]
换用另一个非退化型，只会得到由换基共轭而同构的 Lie 代数。

\ProseParagraphBreak

这里保持的是复双线性型，因此使用普通转置 \(X^{\mathsf T}\)。若保持标准 Hermitian 型，相应条件则是 \(X^*+X=0\)，其中 \(X^*\) 为共轭转置。例如对 \(n\ge2\)，令 \(F=E_{12}-E_{21}\)，则 \(F^{\mathsf T}=-F\)，而 \((iF)^{\mathsf T}=-iF\)，所以 \(F,iF\) 都属于 \(\mathfrak{so}_n(\mathbb C)\)。但 \((iF)^*=iF\)，故 \(iF\) 不是斜 Hermitian 矩阵。前者对复标量封闭，后者只构成实 Lie 代数；尤其不能把这里的复 \(\mathfrak{so}_3\) 与实旋转 Lie 代数 \(\mathfrak{so}_3(\mathbb R)\) 混同。

\begin{proposition}\label{b5:prop:classical-dimensions}
对整数 \(n\ge1\)，复经典 Lie 代数 \(\mathfrak{sl}_n(\mathbb C)\)、\(\mathfrak{so}_n(\mathbb C)\)、\(\mathfrak{sp}_{2n}(\mathbb C)\) 的维数为
\[
\dim\mathfrak{sl}_n=n^2-1,\qquad
\dim\mathfrak{so}_n=\frac{n(n-1)}2,\qquad
\dim\mathfrak{sp}_{2n}=n(2n+1).
\]
相对于标准辛型，\(\mathfrak{sp}_{2n}(\mathbb C)\) 恰好由如下分块矩阵构成，其中 \(A,B,C\in M_n(\mathbb C)\)：
\[
X=\begin{pmatrix}A&B\\C&-A^{\mathsf T}\end{pmatrix},
\qquad B=B^{\mathsf T},\quad C=C^{\mathsf T}.
\]
\end{proposition}
\begin{proof}
迹零是 \(M_n\) 上一个非零线性条件，所以第一式成立。反对称矩阵的对角元为零，每对 \(i<j\) 留一个自由参数，得到第二式。把四个矩阵块代入 \(X^{\mathsf T}J+JX=0\)，便得陈述中的块条件；自由参数数目为
\(n^2+2\cdot n(n+1)/2=n(2n+1)\)。
\end{proof}

\subsection{维数、中心与单性的范围}
\label{b5:subsec:1304:02}

\begin{lemma}[提取共同权分量]\label{b5:lem:classical-weight-projection}
设有限维复向量空间 \(W\) 上的有限族算子 \(H_1,\ldots,H_r\) 交换且可同时对角化。若子空间 \(U\subseteq W\) 被它们保持，则每个 \(u\in U\) 的各个共同特征空间分量仍属于 \(U\)。
\end{lemma}
\begin{proof}
共同特征值元组只有有限多个。可选系数 \(c_i\)，使不同元组在
\(H=\sum_i c_iH_i\) 下具有不同特征值：需要排除的只是有限多个真超平面。对每个特征值 \(\lambda\)，Lagrange 插值多项式
\[
P_\lambda(t)=\prod_{\mu\ne\lambda}\frac{t-\mu}{\lambda-\mu}
\]
给出相应投影 \(P_\lambda(H)\)。子空间 \(U\) 被 \(H\) 及其多项式保持，故被这些投影保持。
\end{proof}

\begin{proposition}\label{b5:prop:classical-simple}
在复数域上，\(\mathfrak{sl}_n\) 对 \(n\ge2\) 为单 Lie 代数，
\(\mathfrak{sp}_{2n}\) 对 \(n\ge1\) 为单 Lie 代数，
\(\mathfrak{so}_n\) 对 \(n=3\) 或 \(n\ge5\) 为单 Lie 代数。因此这些范围内的中心均为零。
\end{proposition}
\begin{proof}
先考虑特殊线性 Lie 代数。
令 \(\mathfrak h\) 为迹零对角矩阵空间。对 \(H=\operatorname{diag}(t_1,\ldots,t_n)\)，有
\[
[H,E_{ij}]=(t_i-t_j)E_{ij}.
\]
这些非零权彼此不同，零权空间为 \(\mathfrak h\)。给定非零理想 \(I\)，\cref{b5:lem:classical-weight-projection}说明：若 \(I\) 的元素有非零非对角分量，则某个 \(E_{ij}\in I\)；否则 \(I\) 含有非零迹零对角矩阵 \(H\)，其对角元不能全相等，故 \([H,E_{ij}]\ne0\) 也使 \(E_{ij}\in I\)。

令 \(H_{ij}=E_{ii}-E_{jj}\)。由 \([E_{ij},E_{ji}]=H_{ij}\) 得到 \(H_{ij}\in I\)，再由 \([H_{ij},E_{ji}]=-2E_{ji}\) 及特征零，得 \(E_{ji}\in I\)。对 \(k\ne i,j\)，与 \(H_{ij}\) 括号先给出 \(E_{ik},E_{ki},E_{jk},E_{kj}\in I\)；随后
\([E_{ki},E_{i\ell}]=E_{k\ell}\) 给出其余非对角单位。全部迹零对角矩阵由 \([E_{ij},E_{ji}]\) 生成，所以 \(I=\mathfrak{sl}_n\)。当 \(n=2\) 时没有额外的 \(k\)，前三个矩阵已经构成一组基。

再考虑辛 Lie 代数。
用 \(\bar i=n+i\) 表示后半部指标，记
\[
\begin{aligned}
H_i&=E_{ii}-E_{\bar i\bar i},&
A_{ij}&=E_{ij}-E_{\bar j\bar i}\quad(i\ne j),\\
B_{ii}&=E_{i\bar i},&
C_{ii}&=E_{\bar i i},\\
B_{ij}&=E_{i\bar j}+E_{j\bar i},&
C_{ij}&=E_{\bar i j}+E_{\bar j i}\quad(i<j).
\end{aligned}
\]
这些矩阵构成一组基。相对于 \(H=\sum_i t_iH_i\) 的伴随作用，它们的权分别是
\(0,t_i-t_j,2t_i,-2t_i,t_i+t_j,-t_i-t_j\)，所有非零权空间均为一维。给定非零理想 \(I\)，若其某个元素有非零对角分量，投影引理把这个零权分量提取为 \(I\) 中的非零对角元；否则，先提取其中一个非零权向量。各非零权空间为一维，所以除以非零系数后，可在 \(I\) 中取得相应的基权向量。再与整个代数中相反权的基向量取括号，理想性质保证结果仍在 \(I\) 中。相应的三种配对为
\[
[A_{ij},A_{ji}]=H_i-H_j,\quad
[B_{ij},C_{ij}]=H_i+H_j,\quad
[B_{ii},C_{ii}]=H_i,
\]
所以无论如何，\(I\) 总含有一个非零对角元 \(H=\sum_i t_iH_i\)。

选择 \(t_i\ne0\)，由 \([H,B_{ii}]=2t_iB_{ii}\)、
\([H,C_{ii}]=-2t_iC_{ii}\) 得 \(B_{ii},C_{ii},H_i\in I\)。
当 \(n=1\) 时，\(H_1,B_{11},C_{11}\) 已是一组基，故 \(I=\mathfrak{sp}_2\)。以下设 \(n\ge2\)。
对 \(j\ne i\)，与 \(H_i\) 括号得到 \(A_{ij},A_{ji}\in I\)，再得到 \(H_i-H_j\)，故所有 \(H_j\) 均在 \(I\) 中。每个非零权至少在某个 \(H_j\) 上非零，与这个 \(H_j\) 括号便把相应基向量收入 \(I\)。所以 \(I=\mathfrak{sp}_{2n}\)。

在 \(\mathfrak{sp}_4\) 中，可以从一个具体的非零权向量 \(B_{11}\in I\) 看见这次传播。矩阵单位的乘法首先给出
\[
\begin{aligned}
[B_{11},C_{11}]&=H_1,&[H_1,C_{11}]&=-2C_{11},\\
[H_1,A_{12}]&=A_{12},&[H_1,A_{21}]&=-A_{21}.
\end{aligned}
\]
因此 \(H_1,C_{11},A_{12},A_{21}\in I\)。再由
\[
[A_{12},A_{21}]=H_1-H_2
\]
得到 \(H_2\in I\)。最后，
\[
\begin{aligned}
[H_2,B_{22}]&=2B_{22},&[H_2,C_{22}]&=-2C_{22},\\
[H_1,B_{12}]&=B_{12},&[H_1,C_{12}]&=-C_{12}
\end{aligned}
\]
把剩下的四个向量也收入 \(I\)。连同最初的 \(B_{11}\)，这正是全部十个基向量。每一步只需括号的一端已在 \(I\) 中，另一端可以是整个代数中的任意元素，这也是使用理想条件的地方。

最后考虑正交 Lie 代数。
记 \(F_{ij}=E_{ij}-E_{ji}\)，其中 \(i<j\)。考虑阶为 \(2^{n-1}\) 的有限交换群
\[
D=\left\{\,\operatorname{diag}(\varepsilon_1,\ldots,\varepsilon_n)
\mid\varepsilon_i=\pm1,\ \prod_i\varepsilon_i=1\,\right\}.
\]
共轭作用将 \(F_{ij}\) 乘以 \(\varepsilon_i\varepsilon_j\)。每个 \(D\) 的元素都是若干同时翻转两个坐标符号的矩阵之积。记同时翻转第 \(i,j\) 个坐标的矩阵为 \(d_{ij}\)，则对 \(X\in\mathfrak{so}_n\) 有
\[
d_{ij}Xd_{ij}^{-1}=X+2[F_{ij},[F_{ij},X]].
\]
这个等式可在基向量上验证：\((\operatorname{ad}F_{ij})^2\) 在只含 \(i,j\) 中一个指标的基向量上为负恒等算子，在 \(F_{ij}\) 及不含这两个指标的基向量上为零；右侧恰好把前一类基向量变号，其余保持不变。Lie 理想被 \(\operatorname{ad}F_{ij}\) 及其多项式保持，故被每个 \(d_{ij}\) 共轭保持，从而被 \(D\) 共轭保持。

选取偶数个坐标变号，正是为了用这个内在的伴随算子多项式保证每个理想都被保持。一般 Lie 代数自同构虽然把理想送到理想，却未必保持一个指定理想。例如对非零 Lie 代数 \(\mathfrak s\)，\(\mathfrak s\oplus\mathfrak s\) 中交换两因子的自同构就把 \(\mathfrak s\oplus0\) 送到 \(0\oplus\mathfrak s\)。因此不能仅凭所有符号翻转都给出自同构，就把它们全部用于提取理想中的分量。

每个指标对 \(\{i,j\}\) 定义 \(D\) 的一维特征标 \(\chi_{ij}(d)=\varepsilon_i\varepsilon_j\)。对 \(n=3\) 或 \(n\ge5\)，不同二元集合给出不同特征标。事实上，两个这样的特征标相同，只可能是两集合相同，或它们的对称差为全部指标；后一个可能性恰好发生于 \(n=4\) 的互补二元集合。因为 \(|D|=2^{n-1}\) 在 \(\mathbb C\) 中非零，可以除以 \(|D|\) 取平均。相应特征空间的投影对 \(X=\sum_{i<j}x_{ij}F_{ij}\) 给出
\[
\frac1{|D|}\sum_{d=\operatorname{diag}(\varepsilon)\in D}
\varepsilon_i\varepsilon_j\,dXd^{-1}=x_{ij}F_{ij}.
\]
不相同的两个特征标相乘为非平凡特征标，其和为零：取它取值为 \(-1\) 的群元素平移求和即可。这正是有限交换群的一维特征标正交投影，在这里把一个理想中的矩阵分离为单个矩阵单位方向。例如在 \(\mathfrak{so}_3\) 中，\(D=\{1,d_{12},d_{13},d_{23}\}\cong C_2\times C_2\)。对 \(X=aF_{12}+bF_{13}+cF_{23}\)，共轭作用分别给出
\[
\begin{aligned}
\operatorname{Ad}_{d_{12}}X&=aF_{12}-bF_{13}-cF_{23},\\
\operatorname{Ad}_{d_{13}}X&=-aF_{12}+bF_{13}-cF_{23},\\
\operatorname{Ad}_{d_{23}}X&=-aF_{12}-bF_{13}+cF_{23}.
\end{aligned}
\]
因此相应 Fourier 投影为
\[
\begin{aligned}
\Pi_{12}(X)&=\frac14\bigl(X+\operatorname{Ad}_{d_{12}}X
-\operatorname{Ad}_{d_{13}}X-\operatorname{Ad}_{d_{23}}X\bigr)\\
&=aF_{12}.
\end{aligned}
\]
若 \(X\in I\) 且 \(a\ne0\)，各共轭作用与它们的线性组合都保持 \(I\)，故 \(F_{12}=a^{-1}\Pi_{12}(X)\in I\)。一般地，某个非零 \(x_{ij}\) 保证理想含有 \(F_{ij}\)。再用
\[
[F_{ij},F_{jk}]=F_{ik}\quad(i,j,k\;\text{互异})
\]
以及反对称记号 \(F_{ji}=-F_{ij}\)，先得到涉及 \(i\) 或 \(j\) 的全部基向量，再得到其余的 \(F_{k\ell}\)。所以非零理想等于 \(\mathfrak{so}_n\)。三类代数均非交换；其中心作为交换理想只能为零。
\end{proof}

上述正交 Lie 代数的证明在 \(n=4\) 时不能分离互补二元集合。例如 \(\chi_{12}=\chi_{34}\)，相应投影只能提取 \(\mathbb CF_{12}\oplus\mathbb CF_{34}\)，不能把这两个方向分开。特征标重合说明这个提取方法在四维时失效，本身还不能证明该代数不单；真正的理想分解需要如下括号计算。\(\mathfrak{so}_2\) 是一维交换代数，中心是自身；\(\mathfrak{so}_1=\mathfrak{sl}_1=0\)。

\begin{proposition}\label{b5:prop:so4-explicit-ideals}
在复 Lie 代数 \(\mathfrak{so}_4\) 中，记 \(F_{ij}=E_{ij}-E_{ji}\)，并令
\[
\begin{aligned}
a_1&=F_{12}+F_{34},&a_2&=F_{13}-F_{24},&a_3&=F_{14}+F_{23},\\
b_1&=F_{12}-F_{34},&b_2&=F_{13}+F_{24},&b_3&=F_{14}-F_{23}.
\end{aligned}
\]
则 \(\mathfrak a=\operatorname{span}_{\mathbb C}\{a_1,a_2,a_3\}\) 与 \(\mathfrak b=\operatorname{span}_{\mathbb C}\{b_1,b_2,b_3\}\) 是两个互相交换的三维单理想，且
\[
\mathfrak{so}_4=\mathfrak a\oplus\mathfrak b.
\]
因而 \(\mathfrak{so}_4\) 不单，但中心为零。
\end{proposition}
\begin{proof}
约定 \(F_{ji}=-F_{ij}\)、\(F_{ii}=0\)。由矩阵单位乘法得到
\[
[F_{ij},F_{k\ell}]
=\delta_{jk}F_{i\ell}-\delta_{ik}F_{j\ell}
-\delta_{j\ell}F_{ik}+\delta_{i\ell}F_{jk}.
\]
把六个向量的定义代入，得到
\[
\begin{aligned}
[a_1,a_2]&=-2a_3,&[a_2,a_3]&=-2a_1,&[a_3,a_1]&=-2a_2,\\
[b_1,b_2]&=2b_3,&[b_2,b_3]&=2b_1,&[b_3,b_1]&=2b_2,
\end{aligned}
\]
以及 \([a_i,b_j]=0\) 对全部 \(1\le i,j\le3\) 成立。例如
\[
\begin{aligned}
[a_1,a_2]&=-F_{23}-F_{14}-F_{14}-F_{23}=-2a_3,\\
[a_1,b_2]&=-F_{23}+F_{14}-F_{14}+F_{23}=0.
\end{aligned}
\]
六个向量由三对可逆的和差变换从 \(F_{ij}\) 得到，故构成一组基。上述括号使两个三维子空间各自闭合且彼此交换，所以它们是理想，并给出所述直和。基 \(a_1/2,a_2/2,a_3/2\) 与 \(-b_1/2,-b_2/2,-b_3/2\) 的循环括号都与 \(\mathfrak{so}_3\) 中 \(F_{12},F_{13},F_{23}\) 的相同，因此两个理想均同构于 \(\mathfrak{so}_3\)，由\cref{b5:prop:classical-simple}均为单 Lie 代数。它们的中心都为零，直和的中心也为零。
\end{proof}

四维情形还说明奇数次符号翻转为什么不能用于上述理想分量的投影。取 \(d=\operatorname{diag}(-1,1,1,1)\)，则 \(d\) 的共轭作用仍是 \(\mathfrak{so}_4\) 的自同构，但
\[
\operatorname{Ad}_d(a_i)=-b_i,\qquad
\operatorname{Ad}_d(b_i)=-a_i\quad(1\le i\le3).
\]
所以它交换 \(\mathfrak a\) 与 \(\mathfrak b\)，而不保持其中任意一个。扩大到所有符号翻转后，虽然可以区分 \(F_{12}\) 与 \(F_{34}\)，相应平均却未必留在原来的理想中。

\subsection{小秩同构与例外情形}
\label{b5:subsec:1304:03}

四维空间的二次外幂是六维的，两个二向量的楔积又落在一维最高外幂中，因而体积元可以把这个楔积解释为一个数。若原空间还带有辛型，其二次外幂中会多出一条由辛型决定的平凡直线；下面先把相应的五维补空间构造出来。

\begin{proposition}\label{b5:prop:sp4-primitive-wedge}
设 \(V\) 是四维复辛空间，取辛基 \(e_1,e_2,f_1,f_2\)，并令
\[
\operatorname{vol}=e_1\wedge e_2\wedge f_1\wedge f_2,\qquad
B(u,v)\operatorname{vol}=u\wedge v
\quad(u,v\in\bigwedge\nolimits^2V).
\]
则 \(B\) 是 \(\bigwedge^2V\) 上非退化的对称双线性型。令
\[
\Omega=e_1\wedge f_1+e_2\wedge f_2,\qquad W=\Omega^\perp.
\]
在 \(\mathfrak{sp}(V)\) 的外幂表示下，\(\mathbb C\Omega\) 是非退化的一维平凡子模，\(W\) 是五维非退化子模，并有子模直和
\[
\bigwedge\nolimits^2V=\mathbb C\Omega\oplus W.
\]
其中 \(W\) 的一组基为
\[
\begin{aligned}
w_1&=e_1\wedge e_2,&w_2&=e_1\wedge f_2,&w_3&=f_1\wedge e_2,\\
w_4&=f_1\wedge f_2,&w_5&=e_1\wedge f_1-e_2\wedge f_2.
\end{aligned}
\]
相对此辛基，令 \(H_i=E_{ii}-E_{2+i,2+i}\)。对 \(H=t_1H_1+t_2H_2\)，上述五个基向量的特征值依次为
\[
t_1+t_2,\quad t_1-t_2,\quad -t_1+t_2,\quad -t_1-t_2,\quad0.
\]
\end{proposition}
\begin{proof}
两个二向量满足 \(u\wedge v=(-1)^{2\cdot2}v\wedge u=v\wedge u\)，故 \(B\) 对称。对辛分块矩阵
\(X=\begin{psmallmatrix}A&B_0\\C&-A^{\mathsf T}\end{psmallmatrix}\)，其中 \(B_0,C\) 对称，直接作用于 \(\Omega\) 得
\[
\begin{aligned}
X\Omega={}&\sum_{i,j=1}^2
\bigl(A_{ji}e_j\wedge f_i-A_{ij}e_i\wedge f_j\bigr)\\
&+\sum_{i,j=1}^2 C_{ji}f_j\wedge f_i
+\sum_{i,j=1}^2 (B_0)_{ji}e_i\wedge e_j=0.
\end{aligned}
\]
第一行交换 \(i,j\) 后成对相消，第二行的两个和分别因对称性与楔积的反对称性相消。因此 \(\mathbb C\Omega\) 是平凡子模。

又有
\[
\Omega\wedge\Omega
=2e_1\wedge f_1\wedge e_2\wedge f_2=-2\operatorname{vol},
\qquad B(\Omega,\Omega)=-2\ne0.
\]
所以 \(B(\Omega,-)\) 是非零泛函，其核 \(W\) 为五维，并且这条直线与 \(W\) 的交为零。对前四个 \(w_i\)，与 \(\Omega\) 楔积时各项都有重复的基向量，故全为零。写 \(u=e_1\wedge f_1\)、\(v=e_2\wedge f_2\)，则
\[
w_5\wedge\Omega=(u-v)\wedge(u+v)
=u\wedge v-v\wedge u=0.
\]
因此五个 \(w_i\) 均在 \(W\) 中。它们线性无关，故为 \(W\) 的一组基。直接按体积元的次序计算配对，得到
\[
B(w_1,w_4)=1,\qquad B(w_2,w_3)=-1,\qquad B(w_5,w_5)=2,
\]
其余基向量配对除对称项外均为零。因此 \(B|_W\) 的 Gram 矩阵在基 \(w_1,w_4,w_2,w_3,w_5\) 下由两个可逆的二阶配对块与非零的一阶块组成，故非退化。这也说明 \(B\) 在整个直和 \(\mathbb C\Omega\oplus W=\bigwedge^2V\) 上非退化。

每个 \(X\in\mathfrak{sp}(V)\) 都迹零。由\cref{b5:prop:lie-top-exterior-trace}，它在 \(\bigwedge^4V\) 上的作用为零，所以外幂作用的 Leibniz 公式给出
\[
B(Xu,v)+B(u,Xv)=0\quad(u,v\in\bigwedge\nolimits^2V).
\]
对 \(w\in W\)，于是 \(B(Xw,\Omega)=-B(w,X\Omega)=0\)，故 \(W\) 被保持，上面的向量空间直和确为子模直和。最后，\(He_i=t_ie_i\)、\(Hf_i=-t_if_i\)，将两因子的特征值相加就得到所列五个值。
\end{proof}

\begin{proposition}[小维数同构]\label{b5:prop:classical-low-rank-isomorphisms}
复数域上有 Lie 代数同构
\[
\mathfrak{sp}_2\cong\mathfrak{sl}_2\cong\mathfrak{so}_3,\qquad
\mathfrak{so}_4\cong\mathfrak{sl}_2\oplus\mathfrak{sl}_2,
\qquad
\mathfrak{sp}_4\cong\mathfrak{so}_5,\qquad
\mathfrak{sl}_4\cong\mathfrak{so}_6.
\]
\end{proposition}
\begin{proof}
二维辛矩阵的块条件恰好是迹零，故 \(\mathfrak{sp}_2=\mathfrak{sl}_2\) 在标准基下直接相等。\(\mathfrak{sl}_2\) 的伴随表示保持非退化对称 Killing 型；其核是零中心。因此它给出单射
\(\mathfrak{sl}_2\rightarrow\mathfrak{so}(\mathfrak{sl}_2,\kappa)\)，两边维数均为 \(3\)，便得到与 \(\mathfrak{so}_3\) 的同构。

设 \(U,W\) 都是带非退化交替型的二维复空间。两个交替型的张量积是 \(U\otimes W\) 上的非退化对称型，且
\[
(A,B)\longmapsto A\otimes I+I\otimes B
\]
把 \(\mathfrak{sl}(U)\oplus\mathfrak{sl}(W)\) 映入相应的四维正交代数。这里通过自然同构 \(\operatorname{End}(U\otimes W)\cong\operatorname{End}(U)\otimes\operatorname{End}(W)\)，定义对第二因子的部分迹为
\[
1\otimes\operatorname{tr}:\operatorname{End}(U)\otimes\operatorname{End}(W)
\longrightarrow\operatorname{End}(U),\qquad
C\otimes D\longmapsto(\operatorname{tr}D)C.
\]
若像为零，应用这个线性映射得到 \(2A+(\operatorname{tr}B)I_U=2A=0\)。同理，对第一因子取部分迹得到 \((\operatorname{tr}A)I_W+2B=2B=0\)。因此映射单射。两边维数为 \(6\)，故得到 \(\mathfrak{so}_4\) 的分解。

在四维复空间 \(V\) 上取基 \(e_1,e_2,e_3,e_4\)，并令 \(\operatorname{vol}=e_1\wedge e_2\wedge e_3\wedge e_4\)，以
\[
B(u,v)\operatorname{vol}=u\wedge v,\qquad u,v\in\bigwedge\nolimits^2V
\]
定义六维空间上的型。两个二向量交换时符号为 \((-1)^{2\cdot2}=1\)，故这个型对称。记 \(e_{ij}=e_i\wedge e_j\)，则
\[
e_{12}\wedge e_{34}=\operatorname{vol},\qquad
e_{13}\wedge e_{24}=-\operatorname{vol},\qquad
e_{14}\wedge e_{23}=\operatorname{vol}.
\]
除了这三对及其交换次序，两个基二向量的楔积都有重复指标而为零。因此在基
\[
e_{12},e_{34},e_{13},-e_{24},e_{14},e_{23}
\]
下，Gram 矩阵由三个 \(\begin{psmallmatrix}0&1\\1&0\end{psmallmatrix}\) 对角块组成，故 \(B\) 非退化。对 \(X\in\mathfrak{sl}(V)\)，\cref{b5:prop:lie-top-exterior-trace}给出 \(X\operatorname{vol}=0\)，于是
\[
\bigl(B(Xu,v)+B(u,Xv)\bigr)\operatorname{vol}
=X(u\wedge v)=0.
\]
所以 \(\mathfrak{sl}(V)\) 保持 \(B\)。对于迹零对角矩阵 \(H=\operatorname{diag}(t_1,t_2,t_3,t_4)\)，有 \(He_{ij}=(t_i+t_j)e_{ij}\)；互补的一对指标满足 \(t_i+t_j=-(t_k+t_\ell)\)，所以三对非零配对方向具有相反的权。这是一般不变性在对角作用中的具体表现。这个表示非零，例如
\(\operatorname{diag}(1,0,-1,0)\) 在 \(e_1\wedge e_2\) 上非零。由\cref{b5:prop:classical-simple}，单 Lie 代数 \(\mathfrak{sl}_4\) 的非零表示核为零，故得到单射 \(\mathfrak{sl}_4\rightarrow\mathfrak{so}_6\)。两边均为 \(15\) 维，因而同构。

最后令四维 \(V\) 带辛型。\cref{b5:prop:sp4-primitive-wedge}说明，同一个六维外幂空间作为 \(\mathfrak{sp}_4\)-表示分解为一维平凡子模 \(\mathbb C\Omega\) 与五维非退化正交子模 \(W\)。因此得到 Lie 代数同态
\[
\mathfrak{sp}_4\longrightarrow\mathfrak{so}(W,B|_W).
\]
这个表示非零：上述 \(w_1=e_1\wedge e_2\in W\)，而 \(H_1w_1=w_1\)。辛代数的单性使其核为零；源与 \(\mathfrak{so}(W,B|_W)\) 均为 \(10\) 维，所以同态为同构，得到 \(\mathfrak{sp}_4\cong\mathfrak{so}_5\)。这里的五维空间是辛型决定的平凡直线的正交补，整个六维外幂表示则是 \(\mathbf1\oplus W\)，其中 \(\mathbf1\) 表示一维平凡表示。
\end{proof}

这些同构是 Lie 代数层面的同构。相应矩阵群还可能相差有限中心或覆盖，不能把本命题直接改写成同名群的同构。后续讨论自旋群和代数群时，会进一步区分这两种结论。
